% =============================================================
% บทที่ 5 - บทสรุป
% =============================================================

% -----------------------------------------------------------
\section{สรุปผลการดำเนินโครงงาน}
% -----------------------------------------------------------

โครงงานนี้มีวัตถุประสงค์เพื่อพัฒนาเว็บแอปพลิเคชันสำหรับระบบเช่าอ่านหนังสือการ์ตูน
พร้อมทั้งออกแบบโมเดลระบบแนะนำหนังสือที่ใช้พฤติกรรมผู้ใช้หลายรูปแบบ และประเมินเปรียบเทียบกับระบบแนะนำแบบพฤติกรรมเดียว
โดยแนวคิดของระบบแนะนำในโครงงานนี้อาศัยหลักการของ Multi-Behavior Recommender System
ซึ่งนำพฤติกรรมของผู้ใช้หลายรูปแบบมาวิเคราะห์ร่วมกันตามห่วงโซ่ View $\rightarrow$ Cart $\rightarrow$ Rent ได้แก่ การคลิกดูรายละเอียด การเพิ่มเข้าตะกร้า และการเช่าหนังสือ ที่เป็นพฤติกรรมเป้าหมาย
ผลการดำเนินงานสรุปเทียบกับวัตถุประสงค์ในบทที่ 1 ได้ดังตารางที่ \ref{tab:obj_result}

\begin{table}[H]
\centering
\caption{สรุปผลเทียบกับวัตถุประสงค์}
\label{tab:obj_result}
\small
\begin{tabular}{|p{0.6cm}|p{3.3cm}|p{5.9cm}|p{1.5cm}|}
\hline
\textbf{ข้อ} & \textbf{วัตถุประสงค์} & \textbf{ผลที่ได้} & \textbf{สถานะ} \\
\hline
1 & เพื่อพัฒนาและสร้างเว็บแอปพลิเคชันต้นแบบระบบแนะนำหนังสือการ์ตูน &
พัฒนาเว็บแอปพลิเคชันด้วย React 19 (Vite) + Tailwind CSS ฝั่ง Frontend และ Django 4.2+ REST Framework + PostgreSQL ฝั่ง Backend ติดตั้งด้วย Docker Compose และนำขึ้นใช้งานบน Railway มีหน้าเว็บ 23 หน้าและ API 44 รายการ ครอบคลุมการสมัครสมาชิกและเข้าสู่ระบบด้วย JWT การเลือกหนังสือที่ชอบ การค้นหา ตะกร้า การจองเช่า การยกเลิก การให้คะแนนและรีวิว (แสดงสถานะที่ปุ่มและป้องกันการแก้ไข) หน้าหนังสือยอดนิยม การเช่าหน้าร้าน การคืนและค่าปรับ Dashboard ผู้ดูแลระบบ และหน้าฝึกและสลับรุ่นโมเดล มีชุดทดสอบอัตโนมัติครอบคลุม 8 test classes รวม 26 test cases และใช้ NotificationModal แทน window.alert เพื่อความสอดคล้องกับธีมเว็บ & บรรลุ \\ \hline
2 & เพื่อออกแบบและนำเสนอโมเดลระบบแนะนำที่ประมวลผลพฤติกรรมผู้ใช้ที่หลากหลาย &
พัฒนาโมเดล MB-CGCN 2 รุ่น ได้แก่ รุ่น v1 (Cart $\rightarrow$ Rent) ฝึกจากชุดข้อมูล MyAnimeList และรุ่น v2 (Click $\rightarrow$ Cart $\rightarrow$ Rent) ที่ใช้น้ำหนักพฤติกรรมแบบเรียนรู้ได้ ฝึกจากข้อมูลพฤติกรรมจำลอง 1,000 ผู้ใช้กับหนังสือ 100 เรื่อง ทั้งสองรุ่นเชื่อมกับหน้า For You ผ่าน RecommenderService พร้อมการแก้ปัญหา Cold-Start ด้วยหนังสือที่ชอบ 4 เรื่อง และการจัดอันดับใหม่ตามแนวหนังสือ & บรรลุ \\ \hline
3 & เพื่อประเมินและเปรียบเทียบประสิทธิภาพกับระบบแนะนำแบบพฤติกรรมเดียวด้วย Recall@K และ NDCG@K &
รุ่น v1 ได้ Recall@20 = 0.0799 และ NDCG@20 = 0.0510 สูงกว่า MF ที่ได้ 0.0584, 0.0400 และ LightGCN แบบพฤติกรรมเดียวที่ได้ 0.0653, 0.0435 บนชุดข้อมูล MyAnimeList รุ่น v2 ได้ Recall@10 = 0.2158 และ NDCG@10 = 0.1318 สูงกว่า MF ที่ได้ 0.1614, 0.0956, LightGCN ที่ได้ 0.1224, 0.0748 และ Popularity ที่ได้ 0.1381, 0.0774 บนชุดข้อมูลจำลอง มีชุดทดสอบอัตโนมัติครอบคลุม Model, API, Permission รวม 26 test cases ทดสอบทั้งการสร้างข้อมูล การเปลี่ยนสถานะ การควบคุมการเข้าถึง และการจัดการตะกร้า & บรรลุ \\
\hline
\end{tabular}
\end{table}

จากตารางที่ \ref{tab:obj_result} โครงงานนี้พัฒนาเว็บแอปพลิเคชันได้ครบตามฟังก์ชันที่วางแผนไว้ โมเดลแบบหลายพฤติกรรมให้ผลการจัดอันดับดีกว่าโมเดลพฤติกรรมเดียวทั้งในชุดข้อมูล MyAnimeList และชุดข้อมูลจำลองภายในระบบ และระบบมีวงจรเก็บพฤติกรรม ฝึกใหม่ และเปิดใช้งานโมเดลได้จากหน้าผู้ดูแลระบบ ซึ่งพร้อมใช้กับข้อมูลของผู้ใช้จริงเมื่อมีการใช้งานมากขึ้น

% -----------------------------------------------------------
\section{ข้อจำกัดของระบบ}
% -----------------------------------------------------------

\subsection{ข้อจำกัดด้านโมเดลและข้อมูล}

\subsubsection{ข้อมูลพฤติกรรมของรุ่น v2 เป็นข้อมูลจำลอง}
ผู้ใช้ 1,000 จาก 1,002 คนในชุดข้อมูลรุ่น v2 เป็นผู้ใช้จำลอง และหนังสือมีเพียง 100 เรื่อง ซึ่งข้อมูลจำลองนี้ถูกสร้างขึ้นเพื่อทดสอบกลไกการฝึกสอนและประเมินประสิทธิภาพของโมเดลเวอร์ชัน 2 ได้ล่วงหน้าก่อนที่ระบบจะมีผู้ใช้งานจริงจำนวนมากพอ ผลการประเมินของรุ่น v2 จึงสะท้อนความสามารถของโมเดลในการเรียนรู้รูปแบบพฤติกรรมแบบ Cascading (Click $\to$ Cart $\to$ Rent) บนข้อมูลจำลองที่ควบคุมได้ แต่ยังไม่ได้ยืนยันกับพฤติกรรมของผู้ใช้จริง และเนื่องจากขนาดและลักษณะของชุดข้อมูลแตกต่างจากโมเดล v1 ซึ่งฝึกจากข้อมูลจริงของ MyAnimeList กว่า 700,000 ผู้ใช้ จึงไม่สามารถใช้เปรียบเทียบตัวเลข Recall และ NDCG ระหว่าง v1 และ v2 โดยตรงได้ ทั้งนี้โครงสร้างและกลไกของโมเดล v2 พร้อมใช้งานกับข้อมูลจริงแล้ว เมื่อระบบสะสมข้อมูลพฤติกรรม CLICK และ RENT อย่างละ 50 รายการขึ้นไปตามเงื่อนไขขั้นต่ำใน retrain\_model\_v2.py ผู้ดูแลระบบจะสามารถฝึกโมเดล v2 ใหม่จากข้อมูลจริงผ่านหน้าแดชบอร์ดได้ทันที

\subsubsection{ผลการแนะนำไม่อัปเดตทันที}
RecommenderService เก็บเวกเตอร์ไว้ในแคชจนกว่าจะเปลี่ยนรุ่นโมเดล พฤติกรรมใหม่ เช่น การเพิ่มลงตะกร้าหรือการคลิก จึงมีผลต่อคำแนะนำก็ต่อเมื่อผู้ดูแลระบบฝึกโมเดลใหม่และกด Activate แล้ว นอกจากนี้ผู้ใช้ที่สมัครหลังการฝึกครั้งล่าสุดจะไม่อยู่ในกราฟ และจะได้คำแนะนำจากหนังสือที่ชอบหรือประวัติการเช่าแทน

\subsubsection{ข้อจำกัดของรุ่น v1 เมื่อใช้งานบนเว็บ}
รุ่น v1 ระบุผู้ใช้ด้วยชื่อผู้ใช้ของ MyAnimeList ผู้ใช้ของเว็บจึงแทบไม่อยู่ในกราฟ และหนังสือของร้านต้องมี mbrs\_id ที่จับคู่ชื่อเรื่องกับ MyAnimeList ได้ หนังสือที่จับคู่ไม่ได้จะไม่ถูกแนะนำโดยโมเดล การฝึกรุ่น v1 ใหม่จากหน้าผู้ดูแลระบบยังต้องใช้ไฟล์ CSV และ script ฝึกสร้างเวกเตอร์ 64 มิติ ขณะที่ Backend โหลดรุ่น v1 ขนาด 32 มิติ ไฟล์ที่ฝึกใหม่จึงจะถูกข้ามเพราะขนาดไม่ตรงกัน

\subsubsection{ความครบถ้วนของการประเมิน}
ไม่มีชุด Validation และรันแต่ละโมเดลเพียงครั้งเดียว รุ่น v1 วัดผลจากการสุ่มผู้ใช้ 1,000 คน รุ่น v2 ฝึกและวัดผลบนชุดข้อมูลจำลองขนาด 1,000 ผู้ใช้ที่แตกต่างจาก v1 จึงไม่ได้เทียบกับ MF และ LightGCN บนข้อมูลชุดเดียวกับ v1 แต่เทียบกับ MF, LightGCN และ Popularity บนข้อมูลจำลองชุดเดียวกันแทน และการวัดผลระหว่างการฝึกของ train\_v2.py ไม่ได้ตัดรายการในชุดฝึกออก 

\subsection{ข้อจำกัดด้านการทำงานของระบบ}

\subsubsection{A/B Testing เชื่อมกับ API แนะนำแล้ว}
ระบบมีตาราง API และหน้าจอสำหรับ A/B Testing โดย RecommendationView เรียก get\_user\_variant เพื่อดึง variant ของผู้ใช้และบันทึกเหตุการณ์ recommendation\_request ผ่าน track\_event พร้อม logging ชื่อ variant และอัลกอริทึมที่ใช้ อย่างไรก็ตาม ระบบยังไม่ได้ใช้ค่า variant.algorithm เพื่อสลับอัลกอริทึมแนะนำจริง และยังไม่มี variant ที่ active ในฐานข้อมูล จึงยังไม่มีผู้ใช้ถูกแบ่งกลุ่มและตัวชี้วัดในหน้าผู้ดูแลระบบเป็น 0

\subsubsection{การบันทึกพฤติกรรมเชื่อมกับ Frontend แล้ว}
พฤติกรรมการเช่า ถูกบันทึกเมื่อผู้ดูแลระบบอนุมัติคำขอ (approve\_order) ผ่าน UserBehaviorLog.objects.create ในลูป items ของคำสั่งเช่า ทำให้บันทึกเฉพาะรายการที่ได้รับอนุมัติจริง การเช่าหน้าร้านไม่ถูกบันทึกเป็นพฤติกรรม โมดูล utils/behaviorLogger.js เชื่อมกับ MangaDetail.jsx และ ComicDetail.jsx แล้ว บันทึก CLICK เมื่อเปิดหน้ารายละเอียดและ ADD\_CART เมื่อเพิ่มลงตะกร้า พร้อมส่ง source จาก location.state (BROWSE, SEARCH, RECOMMENDATION) ส่วน session\_id, duration\_sec และ position ยังไม่ถูกส่ง และยังไม่เชื่อมกับหน้าอื่น (Home, Search, ForYou, Popular)

\subsubsection{งานตามเวลาไม่ทำงานบน Production}
APScheduler เริ่มทำงานเฉพาะเมื่อรันด้วย runserver (เงื่อนไข RUN\_MAIN=true ใน apps.py) บน Railway ที่ใช้ gunicorn จึงไม่มีการเปลี่ยนรายการที่เลยกำหนดเกิน 90 วันเป็น LOST อัตโนมัติ ต้องรันคำสั่ง manage.py update\_lost\_items เอง

\subsubsection{การฝึกโมเดลรันในโปรเซสของเว็บ}
การฝึกเริ่มจาก Thread เบื้องหลังของ Backend ซึ่งใช้ CPU และหน่วยความจำร่วมกับการให้บริการเว็บ (gunicorn 1 Worker) และจำกัดเวลาที่ 1 ชั่วโมง ถ้า Container ถูกรีสตาร์ทระหว่างฝึก รายการ ModelTrainingLog อาจค้างที่สถานะ RUNNING และทำให้สั่งฝึกใหม่ไม่ได้ จนกว่าจะแก้สถานะในฐานข้อมูล

\subsubsection{ข้อจำกัดอื่น}
ระบบไม่มีการชำระเงินออนไลน์และไม่มีการแจ้งเตือนเมื่อสถานะคำสั่งเช่าเปลี่ยน ผู้ใช้เลือกหมายเลขเล่มในหน้ารายละเอียดได้ แต่ Backend เลือกเล่มว่างเล่มแรกเอง API รายการหนังสือและประวัติไม่มีการแบ่งหน้า (Pagination) ผู้ใช้สามารถเลือกหนังสือที่ชอบได้ครั้งเดียวในขั้นตอน onboarding หลังสมัครสมาชิก ระบบบันทึกค่า onboarding\_completed=True และไม่อนุญาตให้แก้ไขภายหลัง

\subsection{ข้อจำกัดด้านฮาร์ดแวร์และซอฟต์แวร์}

\subsubsection{ทรัพยากรประมวลผล}
กราฟรุ่น v1 มีผู้ใช้กว่า 7 แสนคน ต้องฝึกบน GPU ของ Google Colab และไฟล์น้ำหนักกับไฟล์กราฟรวมกันประมาณ 180 MB ซึ่งต้องโหลดไว้ในหน่วยความจำของ Backend ตลอดเวลา จึงต้องจำกัด gunicorn ไว้ที่ 1 Worker และตั้ง OMP\_NUM\_THREADS=1 เพื่อป้องกันหน่วยความจำไม่พอและ Thread ค้าง 

\subsubsection{การตั้งค่าความปลอดภัย}
ค่าเริ่มต้นใน .env.example ใช้ DEBUG=True ซึ่งเปิด CORS ทุก Origin สำหรับการพัฒนา เมื่อนำขึ้นใช้งานจริงต้องตั้ง DEBUG=False, ALLOWED\_HOSTS และ CORS\_ALLOWED\_ORIGINS ให้ถูกต้อง นอกจากนี้โทเค็นถูกเก็บใน localStorage ซึ่งเสี่ยงต่อการโจมตีแบบ XSS มากกว่าการเก็บใน HttpOnly Cookie

% -----------------------------------------------------------
\section{ปัญหาอุปสรรค และแนวทางแก้ไข}
% -----------------------------------------------------------

ในการดำเนินโครงงานนี้พบปัญหาและอุปสรรคบางประการที่ส่งผลต่อกระบวนการพัฒนาระบบ
โดยสามารถแบ่งออกเป็นประเด็นหลักและแนวทางแก้ไขได้ดังตารางที่ \ref{tab:problems}

\begingroup
\small
\begin{longtable}{|p{0.7cm}|p{5.5cm}|p{5.8cm}|}
\caption{ปัญหาอุปสรรคและแนวทางแก้ไข} \label{tab:problems} \\
\hline
\textbf{ลำดับ} & \textbf{ปัญหาที่พบ} & \textbf{แนวทางแก้ไข} \\
\hline
\endfirsthead
\hline
\textbf{ลำดับ} & \textbf{ปัญหาที่พบ} & \textbf{แนวทางแก้ไข} \\
\hline
\endhead
1 & ขาดแคลนข้อมูลพฤติกรรมจริงของผู้ใช้ เนื่องจากเว็บแอปพลิเคชันเพิ่งพัฒนาเสร็จและยังไม่มีผู้ใช้งานจำนวนมาก ทำให้ฝึกโมเดลรุ่น v2 ไม่ได้ &
สร้างข้อมูลจำลองของผู้ใช้ 1,000 คนที่เขียนลงตารางจริงตามกฎของระบบ เพื่อทดสอบวงจรการฝึกและประเมินโมเดล และเตรียมสคริปต์ลบข้อมูลจำลองไว้เมื่อมีข้อมูลจริง \\ \hline
2 & การฝึก MB-CGCN บนกราฟผู้ใช้หลายแสนคนต้องใช้ทรัพยากรสูง  &
ฝึกรุ่น v1 บน Google Colab (GPU Tesla T4) และใช้การสุ่มผู้ใช้ 1,000 คนในการวัดผลระหว่างฝึก ส่วนรุ่น v2 ที่มีขนาดเล็กฝึกบน CPU ใน Container ได้ \\ \hline
3 & Backend บน Railway หยุดทำงานเพราะหน่วยความจำไม่พอและ Worker หมดเวลา เมื่อโหลดโมเดลขนาดใหญ่ใน Worker หลายตัว &
โหลดโมเดลครั้งเดียวในโปรเซสหลักก่อน fork (--preload) ลด Worker เหลือ 1 ตัว ตั้ง timeout 120 วินาที แคชเวกเตอร์ และเพิ่มดัชนี GIN บนคอลัมน์แนวหนังสือ \\ \hline
4 & Connection ของ PostgreSQL ถูกใช้ร่วมกันหลัง fork ทำให้เกิดข้อผิดพลาด SSL connection has been closed unexpectedly &
ปิด Connection ทั้งหมดหลังโหลดโมเดลใน wsgi.py และเปิด CONN\_MAX\_AGE ร่วมกับ CONN\_HEALTH\_CHECKS \\ \hline
5 & การ Build บน Railway ล้มเหลวจากข้อจำกัดการดาวน์โหลด Image ของ Docker Hub (429 Too Many Requests) &
เปลี่ยน Base Image เป็น public.ecr.aws/docker/library \\ \hline
6 & ระบบมีสถานะคำสั่งเช่าหลายขั้น (REQUESTED $\rightarrow$ APPROVED $\rightarrow$ CHECKED\_OUT $\rightarrow$ RETURNED) และสถานะสำเนาที่ต้องเปลี่ยนตาม (AVAILABLE $\rightarrow$ RESERVED $\rightarrow$ RENTED $\rightarrow$ AVAILABLE/LOST/MAINTENANCE) ทำให้ออกแบบให้สอดคล้องกันทุกกรณีได้ยาก และเสี่ยงต่อการจองเล่มเดียวกันซ้อน &
ใช้ Transaction และ select\_for\_update ทุกจุดที่เปลี่ยนสถานะ ตรวจสถานะซ้ำก่อนบันทึก (ตอบ 409 เมื่อไม่ตรง) และใช้ Soft Delete เพื่อรักษาประวัติ \\ \hline
\end{longtable}
\endgroup

% -----------------------------------------------------------
\section{ข้อเสนอแนะในการพัฒนาต่อไป}
% -----------------------------------------------------------

\subsection{การปรับปรุงระบบแนะนำ}

\subsubsection{เก็บข้อมูลจริงและฝึกรุ่น v2 ใหม่}
เปิดให้ผู้ใช้จริงใช้งาน ลบข้อมูลจำลองออก แล้วฝึกรุ่น v2 ใหม่จากหน้า ML Training เมื่อมีข้อมูลพอ จากนั้นเปรียบเทียบกับ MF, LightGCN และ Popularity บนข้อมูลจริงชุดเดียวกัน โดยรันหลายครั้งและทดสอบนัยสำคัญทางสถิติ ผลที่ได้จากข้อมูลจริงจะสะท้อนรูปแบบความสนใจของผู้ใช้งานจริงอย่างแท้จริง และคาดว่าจะให้ประสิทธิภาพที่ดีกว่าการฝึกจากข้อมูลจำลองอีกด้วย

\subsubsection{ปรับการประเมินและการฝึก}
แยกชุด Validation เพื่อเลือก Checkpoint ที่ดีที่สุด แก้ evaluate\_model ของรุ่น v2 ให้ตัดรายการในชุดฝึกออก แก้เงื่อนไขกรองใน prepare\_data.py ให้ตรงกับ Notebook และทำให้ขนาดเวกเตอร์ของรุ่น v1 ตรงกันระหว่างสคริปต์ฝึกและ RecommenderService

\subsubsection{แนะนำตามพฤติกรรมล่าสุด}
หากต้องการให้คำแนะนำเปลี่ยนเร็วขึ้นโดยไม่ต้องฝึกใหม่ทั้งกราฟ สามารถคำนวณเวกเตอร์ผู้ใช้จากค่าเฉลี่ยเวกเตอร์ของหนังสือที่ผู้ใช้คลิกหรือเพิ่มลงตะกร้าล่าสุด แบบเดียวกับวิธี Cold-Start ที่ใช้อยู่ และควรตัดหนังสือที่ผู้ใช้เคยเช่าออกจากผลในทุกกรณี หรือตั้งเวลาฝึกโมเดลใหม่อัตโนมัติ เช่น ทุกสัปดาห์

\subsubsection{เชื่อม A/B Testing กับหน้าแนะนำ}
เรียก ABTestManager.assign\_variant ใน RecommendationView เพื่อเลือกอัลกอริทึมตามกลุ่ม และบันทึก IMPRESSION, CLICK, ADD\_TO\_CART และ RENTAL จากหน้า For You เพื่อวัด CTR และ Conversion Rate ของแต่ละอัลกอริทึมกับผู้ใช้จริง

\subsection{การปรับปรุงฟังก์ชันการทำงาน}

\subsubsection{เพิ่มพฤติกรรมและฟังก์ชันเสริม}
บันทึก source และ position จากหน้า Home, Search และ For You ด้วยโมดูล behaviorLogger.js ที่มีอยู่แล้ว บันทึก RENT เมื่อผู้ดูแลระบบยืนยันการรับหนังสือแทนตอนยื่นคำขอ เพิ่มปุ่มไม่สนใจคำแนะนำ เพื่อเพิ่มสัญญาณพฤติกรรม และเพิ่มการแจ้งเตือนเมื่อคำสั่งเช่าเปลี่ยนสถานะหรือใกล้ครบกำหนดคืน

\subsubsection{ปรับปรุงงานเบื้องหลัง}
ย้ายงานตามเวลา ที่ทำเครื่องหมายหนังสือที่เลยกำหนด ไปเป็น Cron Job ของ Railway หรือ Worker แยก เพื่อให้ทำงานบน Production และย้ายการฝึกโมเดลไปทำงานใน Worker แยก เช่น Celery พร้อมการกู้สถานะ RUNNING ที่ค้าง

\subsubsection{ปรับปรุงประสิทธิภาพ ความปลอดภัย และคุณภาพโค้ด}
เพิ่ม Pagination และ Rate Limiting ให้ /api/mangas/ และ /api/admin/history/ ให้ Backend ใช้หมายเลขเล่มที่ผู้ใช้เลือก ส่วน Popular This Week ใช้ข้อมูลจาก /api/mangas/popular/ พิจารณาเก็บโทเค็นใน HttpOnly Cookie และเพิ่มชุดทดสอบอัตโนมัติของ Django ตาม Test Case ในบทที่ 4
